\sectionheading{5.3 মেশিন ভাষা (Machine Language)}

কম্পিউটারের Processor মানুষের ব্যবহৃত বাংলা, ইংরেজি কিংবা
Python-এর মতো High Level Language সরাসরি বুঝতে পারে না।
Processor মূলত \textbf{0 ও 1-এর নির্দিষ্ট বিন্যাসে তৈরি Binary
Instruction} বুঝে কাজ করে।

কম্পিউটারের সরাসরি বোধগম্য এই Instruction-এর ভাষাই হলো
\textbf{Machine Language}।

%================================================
% Definition
%================================================
\booksubheading{মেশিন ভাষা কাকে বলে?}

\begin{definitionbox}

\textbf{\color{darkTeal}সংজ্ঞা}

যে Programming Language-এ \textbf{0 ও 1-এর Binary Code} ব্যবহার
করে কম্পিউটারকে সরাসরি Instruction দেওয়া হয় এবং যা Processor
সরাসরি বুঝতে ও Execute করতে পারে, তাকে
\textbf{মেশিন ভাষা (Machine Language)} বলে।

\end{definitionbox}

\textbf{সহজভাবে মনে রাখবে—}

\begin{center}

\textbf{
0 ও 1-এর ভাষা
\(\;=\;\)
Machine Language
}

\end{center}

Machine Language কম্পিউটারের Hardware-এর সবচেয়ে কাছাকাছি।
তাই এটি \textbf{Low Level Language}-এর অন্তর্ভুক্ত।

%================================================
% Why Binary
%================================================
\booksubheading{Machine Language-এ 0 ও 1 কেন ব্যবহৃত হয়?}

কম্পিউটারের Electronic Circuit মূলত দুটি অবস্থাকে সহজে প্রকাশ
ও শনাক্ত করতে পারে। যেমন—\textbf{ON/OFF} অথবা
\textbf{High/Low Voltage}।

এই দুটি অবস্থাকে Binary Digit দিয়ে প্রকাশ করা হয়।

\begin{center}

\begin{tabularx}{0.68\linewidth}{|C{0.25\linewidth}|Y|}
\hline

\rowcolor{softBlue}
\textbf{Binary Digit}
&
\multicolumn{1}{c|}{\textbf{সহজ ধারণা}}
\\
\hline

0
&
OFF / Low অবস্থা
\\
\hline

1
&
ON / High অবস্থা
\\
\hline

\end{tabularx}

\end{center}

তাই কম্পিউটারের Instruction শেষ পর্যন্ত 0 ও 1-এর
\textbf{Bit Pattern} আকারে প্রকাশ করা হয়।

%================================================
% Example
%================================================
\begin{examplebox}{Machine Instruction-এর ধারণামূলক উদাহরণ}

ধরা যাক, কোনো Processor-এর একটি Instruction Binary আকারে—

\begin{center}

{\large\textenglish{\ttfamily 10110010 00000101}}

\end{center}

0 ও 1-এর এই নির্দিষ্ট বিন্যাস Processor-এর কাছে একটি নির্দিষ্ট
কাজের নির্দেশ প্রকাশ করতে পারে।

মানুষের পক্ষে এ ধরনের Binary Instruction মনে রাখা ও বোঝা কঠিন,
কিন্তু Processor সরাসরি এটি বুঝতে পারে।

\end{examplebox}

%================================================
% Instruction Structure
%================================================
\booksubheading{Machine Instruction-এর মৌলিক গঠন}

একটি Machine Instruction-এ সাধারণভাবে দুটি গুরুত্বপূর্ণ অংশের
ধারণা পাওয়া যায়—

\begin{center}

\begin{tikzpicture}[node distance=15mm]

\node[databox] (opcode)
{\textbf{Opcode}\\কী কাজ হবে};

\node[processbox, right=of opcode] (operand)
{\textbf{Operand}\\কোন Data/Address-এর ওপর};

\draw[arrow] (opcode) -- (operand);

\end{tikzpicture}

\end{center}

\begin{itemize}

\item \textbf{Opcode (Operation Code):}
Processor-কে কী Operation বা কাজ করতে হবে তা নির্দেশ করে।

\item \textbf{Operand:}
কোন Data, Register বা Memory Location-এর ওপর Operationটি হবে
সে সম্পর্কিত তথ্য প্রকাশ করে।

\end{itemize}

\begin{examplebox}{Opcode ও Operand সহজভাবে বোঝো}

ধরা যাক নির্দেশটি হলো—

\begin{center}

\textbf{A-এর সঙ্গে B যোগ করো}

\end{center}

এখানে \textbf{``যোগ করো''} অংশটি কী Operation করতে হবে তা বোঝায়।
অন্যদিকে \textbf{A ও B} হলো যে Data-এর ওপর কাজটি করা হবে।

Machine Language-এ এই তথ্যগুলো Binary Code-এর মাধ্যমে প্রকাশিত হয়।

\end{examplebox}

%================================================
% How it works
%================================================
\booksubheading{Machine Language কীভাবে কাজ করে?}

Machine Language-এর Instruction সরাসরি Processor-এর
Instruction Set-এর সঙ্গে সম্পর্কিত। তাই Processor Machine Code
সরাসরি Decode করে নির্ধারিত কাজ Execute করতে পারে।

\begin{center}

\begin{tikzpicture}[node distance=13mm]

\node[databox] (code)
{\textbf{Machine Code}\\0 ও 1};

\node[processbox, right=of code] (cpu)
{\textbf{Processor}\\Decode ও Execute};

\node[databox, right=of cpu] (result)
{\textbf{Result}\\কাজ সম্পন্ন};

\draw[arrow] (code) -- (cpu);
\draw[arrow] (cpu) -- (result);

\end{tikzpicture}

\end{center}

\begin{rememberbox}

\textbf{\color{deepOrange}মনে রাখবে}

Machine Language কম্পিউটারের সরাসরি বোধগম্য ভাষা।
তাই Machine Language-এ লেখা Instruction-কে আবার Machine Code-এ
রূপান্তর করার প্রয়োজন হয় না।

\end{rememberbox}

%================================================
% Features
%================================================
\booksubheading{Machine Language-এর বৈশিষ্ট্য}

\begin{tabularx}{\linewidth}{|>{\bfseries}p{0.27\linewidth}|Y|}
\hline

\rowcolor{softBlue}
\multicolumn{1}{|c|}{\textbf{বৈশিষ্ট্য}}
&
\multicolumn{1}{c|}{\textbf{ব্যাখ্যা}}
\\
\hline

Binary ভিত্তিক
&
Instruction 0 ও 1-এর Binary Code দিয়ে লেখা হয়।
\\
\hline

সরাসরি বোধগম্য
&
Processor Machine Instruction সরাসরি বুঝতে ও Execute করতে পারে।
\\
\hline

Low Level Language
&
Hardware-এর খুব কাছাকাছি অবস্থান করে।
\\
\hline

Machine Dependent
&
নির্দিষ্ট Processor বা Machine-এর Instruction Set-এর ওপর নির্ভরশীল।
\\
\hline

Translator প্রয়োজন নেই
&
Machine Code Processor সরাসরি Execute করতে পারে।
\\
\hline

Program লেখা কঠিন
&
Binary Instruction মনে রাখা ও হাতে লেখা মানুষের জন্য কঠিন।
\\
\hline

Error শনাক্ত কঠিন
&
দীর্ঘ 0 ও 1-এর মধ্যে ভুল খুঁজে বের করা সময়সাপেক্ষ।
\\
\hline

\end{tabularx}

%================================================
% Machine Dependence
%================================================
\booksubheading{Machine Language কেন Machine Dependent?}

প্রতিটি Processor বা Processor Family-এর নিজস্ব
\textbf{Instruction Set} থাকতে পারে।

অর্থাৎ একটি নির্দিষ্ট Binary Pattern কোনো Processor-এর কাছে
যে Instruction প্রকাশ করে, অন্য ধরনের Processor-এর কাছে একই
Binary Pattern একই অর্থ নাও প্রকাশ করতে পারে।

তাই একটি Machine বা Processor-এর জন্য লেখা Machine Language
Program অন্য ধরনের Machine-এ সরাসরি ব্যবহার করা সবসময় সম্ভব নয়।

\begin{examplebox}{Machine Dependence সহজভাবে বোঝো}

ধরা যাক, \textbf{Computer-A} ও \textbf{Computer-B}-তে দুই ধরনের
Processor রয়েছে।

Computer-A-এর জন্য লেখা Binary Instruction যদি Computer-B-এর
Instruction Set-এর সঙ্গে না মেলে, তবে একই Machine Program
Computer-B-তে সরাসরি কাজ করবে না।

এই কারণে Machine Language-কে
\textbf{Machine Dependent Language} বলা হয়।

\end{examplebox}

%================================================
% Advantages
%================================================
\booksubheading{Machine Language-এর সুবিধা}

\begin{itemize}

\item \textbf{Processor সরাসরি বুঝতে পারে—}
Machine Language-এর Instruction সরাসরি Execute করা যায়।

\item \textbf{অনুবাদের অতিরিক্ত ধাপ প্রয়োজন হয় না—}
Instruction ইতোমধ্যে Machine-এর বোধগম্য Binary রূপে থাকে।

\item \textbf{Hardware-এর সঙ্গে সরাসরি কাজ করা যায়—}
Processor-এর Instruction Set-এর ওপর সরাসরি নিয়ন্ত্রণ পাওয়া যায়।

\item \textbf{Memory ও Hardware-level Operation নির্দিষ্টভাবে নিয়ন্ত্রণ
করা সম্ভব।}

\end{itemize}

%================================================
% Disadvantages
%================================================
\booksubheading{Machine Language-এর অসুবিধা}

Machine Language কম্পিউটারের জন্য সরাসরি বোধগম্য হলেও মানুষের
জন্য এটি ব্যবহার করা খুবই কঠিন।

\begin{compacttwocolpoints}

\item 0 ও 1 ব্যবহার করে Program লেখা অত্যন্ত কঠিন।

\item Binary Instruction মনে রাখা কঠিন।

\item Program লেখা সময়সাপেক্ষ।

\item Program পড়া ও Logic বুঝতে অসুবিধা হয়।

\item একটি Bit-এর ভুলেও Instruction বা Result পরিবর্তিত হতে পারে।

\item Error শনাক্ত ও সংশোধন করা কঠিন।

\item Program পরিবর্তন ও Maintenance করা জটিল।

\item Machine Dependent হওয়ায় Portability কম।

\item বড় Program Machine Language-এ লেখা অত্যন্ত কষ্টসাধ্য।

\end{compacttwocolpoints}

%================================================
% Comparison
%================================================
\booksubheading{Machine Language-এর সুবিধা ও অসুবিধা এক নজরে}

\begin{tabularx}{\linewidth}{|Y|Y|}
\hline

\rowcolor{softBlue}
\multicolumn{1}{|c|}{\textbf{সুবিধা}}
&
\multicolumn{1}{c|}{\textbf{অসুবিধা}}
\\
\hline

Processor সরাসরি বুঝতে পারে
&
মানুষের জন্য লেখা ও বোঝা কঠিন
\\
\hline

অতিরিক্ত Translator প্রয়োজন নেই
&
Error খুঁজে বের করা কঠিন
\\
\hline

Hardware-এর সঙ্গে সরাসরি কাজ করা যায়
&
Machine Dependent
\\
\hline

Instruction সরাসরি Execute করা যায়
&
Program তৈরি ও পরিবর্তন সময়সাপেক্ষ
\\
\hline

\end{tabularx}

%================================================
% Generation
%================================================
\booksubheading{Machine Language ও First Generation Language}

Programming Language-এর Generation অনুযায়ী Machine Language-কে
\textbf{First Generation Language (1GL)} বলা হয়।

\begin{center}

\textbf{
Machine Language
\(\;=\;\)
First Generation Language (1GL)
}

\end{center}

এটি Programming Language-এর সবচেয়ে প্রাথমিক স্তরের ভাষাগুলোর
একটি এবং Hardware-এর সঙ্গে সরাসরি সম্পর্কিত।

%================================================
% Why next language
%================================================
\booksubheading{Machine Language থেকে নতুন ভাষার প্রয়োজন কেন হলো?}

Machine Language-এ Program লেখা, Binary Instruction মনে রাখা,
Program পরিবর্তন করা এবং Error খুঁজে বের করা Programmer-এর জন্য
খুব কঠিন ছিল।

এই সমস্যা কমানোর জন্য 0 ও 1-এর পরিবর্তে সহজ প্রতীকী শব্দ বা
\textbf{Mnemonic Code} ব্যবহার করে Instruction লেখার ধারণা তৈরি হয়।
এর ফলে \textbf{Assembly Language}-এর বিকাশ ঘটে।

\begin{center}

\textbf{
Binary Code
\(\longrightarrow\)
Program লেখা কঠিন
\(\longrightarrow\)
Mnemonic-এর প্রয়োজন
\(\longrightarrow\)
Assembly Language
}

\end{center}

পরবর্তী অংশে Assembly Language বিস্তারিত আলোচনা করা হবে।

%================================================
% Final Reminder
%================================================
\begin{rememberbox}

\textbf{\color{deepOrange}মনে রাখবে}

\begin{itemize}

\item Machine Language হলো \textbf{0 ও 1-এর Binary ভাষা}।

\item এটি একটি \textbf{Low Level Language}।

\item Machine Language হলো
\textbf{First Generation Language (1GL)}।

\item Processor Machine Language
\textbf{সরাসরি বুঝতে ও Execute করতে পারে}।

\item Machine Instruction-এ
\textbf{Opcode ও Operand}-এর ধারণা গুরুত্বপূর্ণ।

\item Machine Language হলো
\textbf{Machine Dependent}।

\item Machine Language-এর জন্য Machine Code তৈরির
Translator প্রয়োজন হয় না।

\item কম্পিউটারের জন্য সরাসরি বোধগম্য হলেও মানুষের জন্য
Program লেখা, পড়া ও Debug করা কঠিন।

\end{itemize}

\end{rememberbox}

\small
